%==========================================================================
%  TRES0312 Fysikk — Foiler-mal
%
%  MØrk/lys modus: kommenter ut én av de to linjene nedenfor.
%==========================================================================
\newif\ifdarkmode
\darkmodetrue      % ← mørk modus
%\darkmodefalse   % ← lys modus  (fjern % foran og legg % foran linjen over)

\documentclass[aspectratio=169, 12pt]{beamer}

\usepackage{amd-beamer}
\usetikzlibrary{overlay-beamer-styles}
\usetikzlibrary{calc}

%--------------------------------------------------------------------------
%  DOKUMENTINFO
%--------------------------------------------------------------------------
\title{Temperatur, varme og faseoverganger}
\subtitle{TRES0312 --- kapittel 8.1 og 8.2}
\author{Ben David Normann}
\date{\today}

%--------------------------------------------------------------------------
%  SEKSJONSKONTEKST
%--------------------------------------------------------------------------
\setcounter{section}{8}   % Kapittel 8: "Termofysikk"

\begin{document}

%--------------------------------------------------------------------------
%  SLIDE 1 — TITTELSIDE
%--------------------------------------------------------------------------
\begin{frame}[plain]
  \titlepage
  \dekkerdelkap{kapittel 8.1 og 8.2}
  \vspace{0.3cm}
  \begin{center}
  {\bfseries\color{repcol} Første time: indre energi, temperatur og varme. Etter pausen: faseoverganger.}
  \end{center}
\end{frame}

%--------------------------------------------------------------------------
%  SLIDE 2 — Tre kjappe fra forrige gang
%--------------------------------------------------------------------------
\begin{frame}{}

  \repetisjon{Tre kjappe fra forrige gang}{%
    \begin{enumerate}[1.]
      \item Et kast: Hvilken bevegelse har vi i $x$-retning, og hvilken i $y$-retning?
      \item Når er den mekaniske energien $E_\text{K} + E_\text{P}$ bevart?
      \item Hva er trykket i dybden $h$ i en væske i ro?
    \end{enumerate}
  }

\end{frame}

%==========================================================================
%  DEL 1: INDRE ENERGI, TEMPERATUR OG VARME (8.1)
%==========================================================================

%--------------------------------------------------------------------------
%  SLIDE 3 — Til ettertanke: havet
%--------------------------------------------------------------------------
\begin{frame}{Indre energi}

  \tanke{}{%
  Er det så farlig om temperaturen i havet stiger med en grad eller to?
  }
  \pause
  \svar{
  Ja. Temperatur handler om \textit{energi}. Når hele havet blir en grad varmere, har det fått tilført enormt mye energi. Den energien påvirker vær, strømmer, is og liv i havet.
  }

\end{frame}

%--------------------------------------------------------------------------
%  SLIDE 4 — Indre energi
%--------------------------------------------------------------------------
\begin{frame}[fragile]{Indre energi}

  Alt stoff består av molekyler i stadig bevegelse. Den samlede energien til molekylene kalles \textbf{indre energi} $U$.

  \begin{center}
\resizebox{0.92\linewidth}{!}{%
\begin{tikzpicture}[
  font=\small,
  every node/.style={text=white},
  >={Stealth[length=5pt,width=4pt]},
  pil/.style={->, thick, white!55},
  kat/.style={draw=#1!80, fill=#1!35!black, rounded corners=4pt, thick,
              minimum width=3.4cm, minimum height=0.8cm, font=\small\bfseries},
  leaf/.style={draw=#1!70, fill=#1!12!black, rounded corners=4pt, thick,
               minimum width=3.2cm, minimum height=2.5cm},
  atom/.style={circle, shading=ball, ball color=#1, minimum size=0.36cm, inner sep=0pt},
  merk/.style={font=\scriptsize, text=#1!40!white, anchor=south},
]
  % Nivå 1
  \node[draw=headCol!80, fill=headCol!45!black, rounded corners=4pt, thick,
        minimum width=4.2cm, minimum height=0.9cm, font=\bfseries] (U) at (0,0) {Indre energi $U$};
  % Nivå 2
  \node[kat=defscol] (K) at (-3.8,-1.6) {Kinetisk};
  \node[kat=red]     (P) at ( 5.3,-1.6) {Potensiell};
  \draw[pil] (U.south) -- ++(0,-0.35) -| (K.north);
  \draw[pil] (U.south) -- ++(0,-0.35) -| (P.north);

  % Nivå 3: bokser
  \foreach \x/\n/\c in {-7.2/tr/defscol, -3.8/vi/defscol, -0.4/ro/defscol, 3.6/fa/red, 7.0/ke/red} {
    \node[leaf=\c] (\n) at (\x,-4.0) {};
  }
  \foreach \n in {tr,vi,ro} { \draw[pil] (K.south) -- ++(0,-0.3) -| (\n.north); }
  \foreach \n in {fa,ke} { \draw[pil] (P.south) -- ++(0,-0.3) -| (\n.north); }

  % Translasjon
  \node[font=\small\bfseries, anchor=north] at ($(tr.north)+(0,-0.12)$) {Translasjon};
  \draw[white!45, dashed, thick] ($(tr.center)+(-1.1,-0.15)$) -- ($(tr.center)+(-0.35,-0.15)$);
  \node[atom=oppgcol!45!white] at ($(tr.center)+(-0.2,-0.15)$) {};
  \draw[->, thick, defscol!60!white] ($(tr.center)+(0.05,-0.15)$) -- ($(tr.center)+(1.0,-0.15)$);

  % Vibrasjon
  \node[font=\small\bfseries, anchor=north] at ($(vi.north)+(0,-0.12)$) {Vibrasjon};
  \draw[decorate, decoration={coil, segment length=3pt, amplitude=2.5pt}, white!70, thick]
       ($(vi.center)+(-0.45,-0.15)$) -- ($(vi.center)+(0.45,-0.15)$);
  \node[atom=oppgcol!45!white] at ($(vi.center)+(-0.6,-0.15)$) {};
  \node[atom=oppgcol!45!white] at ($(vi.center)+(0.6,-0.15)$) {};
  \draw[<->, thick, defscol!60!white] ($(vi.center)+(-0.95,-0.65)$) -- ($(vi.center)+(-0.25,-0.65)$);
  \draw[<->, thick, defscol!60!white] ($(vi.center)+(0.25,-0.65)$) -- ($(vi.center)+(0.95,-0.65)$);

  % Rotasjon
  \node[font=\small\bfseries, anchor=north] at ($(ro.north)+(0,-0.12)$) {Rotasjon};
  \draw[white!70, thick] ($(ro.center)+(-0.45,-0.15)$) -- ($(ro.center)+(0.45,-0.15)$);
  \node[atom=oppgcol!45!white] at ($(ro.center)+(-0.5,-0.15)$) {};
  \node[atom=oppgcol!45!white] at ($(ro.center)+(0.5,-0.15)$) {};
  \draw[->, thick, defscol!60!white] ($(ro.center)+(0.95,0.05)$) arc (15:330:0.98 and 0.5);

  % Bindinger mellom molekylene
  \node[font=\footnotesize\bfseries, anchor=north, align=center] at ($(fa.north)+(0,-0.12)$) {Bindinger mellom\\[-2pt]molekylene};
  \foreach \i in {0,1,2} \foreach \j in {0,1} {
    \coordinate (a\i\j) at ($(fa.center)+(-0.5+0.5*\i,-0.45+0.5*\j)$);
  }
  \foreach \i in {0,1,2} \draw[white!60, thick] (a\i0) -- (a\i1);
  \foreach \j in {0,1} \draw[white!60, thick] (a0\j) -- (a2\j);
  \foreach \i in {0,1,2} \foreach \j in {0,1} { \node[atom=headCol!70!white, minimum size=0.3cm] at (a\i\j) {}; }
  \node[merk=red] at ($(fa.south)+(0,0.08)$) {(faseoverganger)};

  % Kjemiske bindinger
  \node[font=\small\bfseries, anchor=north, align=center] at ($(ke.north)+(0,-0.12)$) {Kjemiske\\[-2pt]bindinger};
  \coordinate (o) at ($(ke.center)+(0,-0.05)$);
  \draw[white!65, thick] (o) -- ++(-0.5,-0.35) (o) -- ++(0.5,-0.35);
  \node[atom=red, minimum size=0.45cm] at (o) {};
  \node[atom=white, minimum size=0.3cm] at ($(o)+(-0.55,-0.38)$) {};
  \node[atom=white, minimum size=0.3cm] at ($(o)+(0.55,-0.38)$) {};
  \node[merk=red] at ($(ke.south)+(0,0.08)$) {(kjemiske reaksjoner)};
\end{tikzpicture}%
}
  \end{center}

  \vspace{-0.6em}
  {\footnotesize Temperaturen $T$ er et mål på bevegelsen til molekylene (den kinetiske delen). Energien som skal til for å bryte bindingene (den potensielle delen), viser seg ikke på temperaturen.\par}

\end{frame}

%--------------------------------------------------------------------------
%  SLIDE 5 — Til ettertanke: lavest mulig temperatur
%--------------------------------------------------------------------------
\begin{frame}{Temperatur}

  \tanke{}{%
  Finnes det en lavest mulig temperatur?
  }
  \pause
  \svar{
  Ja. Jo mindre molekylene beveger seg, desto lavere er temperaturen. Det finnes en temperatur der molekylene ikke kan miste mer bevegelsesenergi: det \textit{absolutte nullpunktet}, $-273.15\,^\circ$C. Den \textit{absolutte temperaturen} $T$ måles fra dette nullpunktet, i kelvin (K).
  }

\end{frame}

%--------------------------------------------------------------------------
%  SLIDE 6 — Definisjon: gjennomsnittlig kinetisk energi
%--------------------------------------------------------------------------
\begin{frame}{Temperatur}

  \defnat{8.1.1}{Gjennomsnittlig kinetisk energi}{
  Den gjennomsnittlige kinetiske energien knyttet til translasjon er for hvert molekyl i en gass
  \[
    \bar{E}_\text{K} = \frac{3}{2}\,k_\text{B}T
  \]
  der $k_\text{B} = 1.38\cdot 10^{-23}\,\mathrm{J/K}$ er Boltzmanns konstant og $T$ er den absolutte temperaturen.
  }
  \pause
  \graabox{$\displaystyle T = t + 273$ \qquad ($T$ i kelvin, $t$ i grader celsius)}
  \pause
  En temperatur\textit{endring} er like stor på begge skalaene: $\Delta T = \Delta t$.

\end{frame}

%--------------------------------------------------------------------------
%  SLIDE 7 — Aktivitet: farten til nitrogenmolekyler
%--------------------------------------------------------------------------
%  Fasit: T = 25 + 273 = 298 K.
%         E_K = 3/2 * 1.38e-23 * 298 = 6.17e-21 J.
%         v = sqrt(2 E_K / m) = sqrt(2*6.17e-21/4.65e-26) = 515 m/s  (ca. 1850 km/t).
\begin{frame}{Aktivitet}

  \aktivitet{Farten til nitrogenmolekyler}{%
  Lufta i rommet holder $25\,^\circ$C. Et nitrogenmolekyl har massen $m = 4.65\cdot 10^{-26}\,\mathrm{kg}$.
  \begin{enumerate}[a)]
    \item Finn den gjennomsnittlige kinetiske energien til et nitrogenmolekyl.
    \item Omtrent hvor fort beveger nitrogenmolekylene seg? \textit{(Hint: $E_\text{K} = \tfrac{1}{2}mv^2$.)}
  \end{enumerate}
  }

\end{frame}

%--------------------------------------------------------------------------
%  SLIDE 8 — Aktivitet: dobbelt så stor energi
%--------------------------------------------------------------------------
%  Fasit: T1 = 17 + 273 = 290 K. Dobbel energi gir dobbel T: T2 = 580 K = 307 °C.
\begin{frame}{Aktivitet}

  \aktivitet{Dobbelt så stor energi}{%
  En gass har temperaturen $17\,^\circ$C. Vi varmer den opp til den gjennomsnittlige kinetiske energien til molekylene er dobbelt så stor. Hva er temperaturen da, i kelvin og i celsius?
  }

  \pause
  \merk{Husk:}{%
  Dobbelt så stor energi betyr dobbelt så høy temperatur \textit{i kelvin}.
  }

\end{frame}

%--------------------------------------------------------------------------
%  SLIDE 9 — Definisjon: varme
%--------------------------------------------------------------------------
\begin{frame}{Varme}

  \defnat{8.1.2}{Varme}{
  Varme $Q$ er indre energi som overføres fra ett system til et annet \textit{som følge av en temperaturforskjell}.
  }
  \pause
  \begin{itemize}
    \item Et system \textit{har} en temperatur, men det «har» ikke varme.
    \item<3-> Varme overføres ved \textbf{varmeledning} (kontakt), \textbf{konveksjon} (strømmende væske eller gass) og \textbf{stråling} (som fra sola).
  \end{itemize}

\end{frame}

%--------------------------------------------------------------------------
%  SLIDE 10 — Til ettertanke: to kjeler
%--------------------------------------------------------------------------
\begin{frame}{Varme}

  \tanke{}{%
  En kjele med $1\,\mathrm{L}$ vann og en kjele med $5\,\mathrm{L}$ vann har begge $60\,^\circ$C. Hvilken har høyest temperatur? Hvilken har mest indre energi? Hvilken avgir mest varme når den kjøles ned til romtemperatur?
  }
  \pause
  \svar{
  Samme temperatur. Men den store kjelen har fem ganger så mange molekyler, og dermed mest indre energi. Den avgir også fem ganger så mye varme når den kjøles ned.
  }

\end{frame}

%--------------------------------------------------------------------------
%  SLIDE 11 — Oppsummering: del 1
%--------------------------------------------------------------------------
\begin{frame}{}

  \oppsum{Indre energi, temperatur og varme}{%
  \begin{enumerate}[1.]
    \item[] Indre energi $U$: kinetisk (bevegelse) og potensiell (bindinger).
    \item[] Absolutt temperatur er et mål på bevegelsen: $\bar{E}_\text{K} = \tfrac{3}{2}k_\text{B}T$, med $T = t + 273$.
    \item[] Varme $Q$ er energi som overføres på grunn av en temperaturforskjell.
  \end{enumerate}
  }

\end{frame}

%--------------------------------------------------------------------------
%  SLIDE 12 — PAUSE
%--------------------------------------------------------------------------
\begin{frame}
\begin{center}
  \Huge{Pause}\\[8pt]
  {\large Etter pausen: faseoverganger}
\end{center}
\end{frame}

%==========================================================================
%  DEL 2: FASEOVERGANGER (8.2)
%==========================================================================

%--------------------------------------------------------------------------
%  SLIDE 13 — Til ettertanke: poteter på full plate
%--------------------------------------------------------------------------
\begin{frame}{Faseoverganger}

  \tanke{}{%
  Potetene koker. Blir de fortere ferdige om du skrur plata opp på full effekt?
  }
  \pause
  \svar{
  Nei. Vannet blir ikke varmere enn $100\,^\circ$C, uansett hvor mye energi du tilfører. Den ekstra energien går bare til å koke bort vannet fortere.
  }

\end{frame}

%--------------------------------------------------------------------------
%  SLIDE 14 — Oppvarmingskurven for vann
%--------------------------------------------------------------------------
\begin{frame}{Faseoverganger}

  \begin{columns}[c]
    \begin{column}{0.44\textwidth}
      \begin{itemize}
        \item<1-> Under en \textbf{faseovergang} er temperaturen \textit{konstant}.
        \item<2-> Energien går til å bryte bindinger mellom molekylene, ikke til mer bevegelse.
        \item<3-> Fordamping krever mye mer energi enn smelting.
      \end{itemize}
    \end{column}
    \begin{column}{0.54\textwidth}
      \centering
      \begin{tikzpicture}[xscale=0.62, yscale=0.62, font=\scriptsize]
        \draw[->, textcolor!70] (0,0) -- (9.0,0) node[right] {$E$/kJ};
        \draw[->, textcolor!70] (0,0) -- (0,7.6) node[above] {$t$/$^\circ$C};
        \foreach \x/\l in {0/0, 2/800, 4/1600, 6/2400, 8/3200} \draw[textcolor!70] (\x,0.08) -- (\x,-0.08) node[below] {\l};
        \foreach \y/\l in {0/{-20}, 1/0, 6/100} \draw[textcolor!70] (0.08,\y) -- (-0.08,\y) node[left] {$\l$};
        \draw[line width=2pt, defscol] (0,0) -- (0.105,1);
        \draw[line width=2pt, oppgcol] (0.105,1) -- (0.94,1);
        \draw[line width=2pt, defscol] (0.94,1) -- (1.99,6);
        \draw[line width=2pt, oppgcol] (1.99,6) -- (7.64,6);
        \draw[line width=2pt, defscol] (7.64,6) -- (8.5,7.2);
        \node[right] at (0.15,0.45) {is};
        \node[above] at (0.52,1.05) {smelting};
        \node[right] at (1.45,3.5) {vann};
        \node[above] at (4.8,6.05) {fordamping, $2259\,\mathrm{kJ}$};
        \node[below] at (0.6,0.95) {$334\,\mathrm{kJ}$};
        \node[left] at (8.05,6.95) {damp};
      \end{tikzpicture}
      \\[-1mm]{\scriptsize Oppvarming av $1\,\mathrm{kg}$ vann fra $-20\,^\circ$C}
    \end{column}
  \end{columns}

\end{frame}

%--------------------------------------------------------------------------
%  SLIDE 15 — Aktivitet: kokende vann på full plate
%--------------------------------------------------------------------------
%  Fasit: a) Energien går til å bryte bindinger (fordamping), T holder seg på 100 °C.
%         b) Fra kurven: 0 -> 100 °C krever ca. 420 kJ, fordamping 2259 kJ.
%            Forhold 2259/420 = 5.4. Tid: 5.4 * 3.0 min = 16 min.
\begin{frame}{Aktivitet}

  \aktivitet{Kokende vann på full plate}{%
  En kjele med $1.0\,\mathrm{kg}$ vann står på en plate som hele tiden gir like mye energi per sekund. Det tar $3.0\,\mathrm{min}$ å varme vannet fra $0\,^\circ$C til $100\,^\circ$C.
  \begin{enumerate}[a)]
    \item Hvorfor blir ikke vannet varmere enn $100\,^\circ$C?
    \item Bruk oppvarmingskurven: Omtrent hvor lang tid tar det før alt vannet har kokt bort?
  \end{enumerate}
  }

\end{frame}

%--------------------------------------------------------------------------
%  SLIDE 16 — Til ettertanke: svette og damp
%--------------------------------------------------------------------------
\begin{frame}{Faseoverganger}

  \tanke{}{%
  Hvorfor blir vi avkjølt når svetten fordamper? Og hvorfor gir ett gram damp på $100\,^\circ$C mer energi til huden enn ett gram vann på $100\,^\circ$C?
  }
  \pause
  \svar{
  Fordamping krever energi, og den tas fra huden. Motsatt: Når damp kondenserer på huden, avgis den samme energien, $2259\,\mathrm{kJ}$ per kilogram, før vannet i tillegg kjøles ned.
  }

\end{frame}

%--------------------------------------------------------------------------
%  SLIDE 17 — Merk: trykk
%--------------------------------------------------------------------------
\begin{frame}{Faseoverganger}

  \merk{Faseoverganger avhenger av trykk}{%
  Høyt til fjells er lufttrykket lavere, og vann koker ved under $100\,^\circ$C. Det tar derfor lenger tid å koke egg på fjellet: vannet blir aldri like varmt.
  }

\end{frame}

%--------------------------------------------------------------------------
%  SLIDE 18 — Oppsummering: del 2
%--------------------------------------------------------------------------
\begin{frame}{}

  \oppsum{Faseoverganger}{%
  \begin{enumerate}[1.]
    \item[] Smelting, størkning, fordamping og kondensering er faseoverganger.
    \item[] Under en faseovergang er temperaturen konstant.
    \item[] For vann: smelting krever $334\,\mathrm{kJ/kg}$, fordamping $2259\,\mathrm{kJ/kg}$.
  \end{enumerate}
  }

\end{frame}

%--------------------------------------------------------------------------
%  SLIDE 19 — Neste gang
%--------------------------------------------------------------------------
\begin{frame}
\begin{center}
  \Huge{Neste gang: kalorimetri og termodynamikkens 1.\ lov}
  \vspace{6mm}
\end{center}
\end{frame}

\end{document}
